\section{Introduction}
\label{sec:intro}

\subsection{The question}

Consider a display that hovers beside a person while they work standing, and
follows them when they move. The engineering question is not whether a
multirotor can carry a screen; multirotors carry cameras of comparable mass
routinely. The question is whether the coupled system closes, and if it does,
whether the machine it produces is one a person would tolerate at arm's length
for half an hour.

The coupling is the whole difficulty. Write it as a cycle:
\begin{equation}
  m \;\longrightarrow\; T \;\longrightarrow\; P \;\longrightarrow\; E
    \;\longrightarrow\; \mbat \;\longrightarrow\; m ,
  \label{eq:cycle}
\end{equation}
gross mass sets the thrust required to hover, thrust sets the power through
momentum theory, power integrated over the mission sets the energy, energy
divided by specific energy sets the battery mass, and the battery is part of
the gross mass. A second cycle runs through the dynamics,
\begin{equation}
  m, J \;\longrightarrow\; \text{plant} \;\longrightarrow\; u(t)
    \;\longrightarrow\; P(t) \;\longrightarrow\; E \;\longrightarrow\; \mbat
    \;\longrightarrow\; m, J ,
  \label{eq:cycle2}
\end{equation}
because the power actually consumed depends on what the controller does, which
depends on the inertia, which depends on the mass distribution, which depends
on the battery.

Neither cycle is a mere iteration to be carried out until convergence. The
first has, under a fixed rotor geometry, a genuine bifurcation: past a
critical payload there is no fixed point at all. That is the mathematical
content of \cref{sec:fold}, and it is the reason a paper is warranted rather
than a spreadsheet.

\subsection{What this paper does}

We construct the problem in three layers and then solve them together.

\paragraph{Layer one: algebra.} Momentum theory over an actuator disk
(\cref{sec:momentum}) gives hover power as a function of thrust and disk area.
A blade-element profile term (\cref{sec:blade}) turns the figure of merit from
an assumption into a consequence of blade geometry. Substituting into a mass
budget (\cref{sec:closure}) produces a single scalar equation whose solvability
is the existence of the design. \Cref{sec:fold} analyses that equation
completely: a trichotomy theorem, the closed form of the critical payload, the
identification of the transition as a saddle-node bifurcation, and a proof that
the naive sizing iteration converges precisely to the physically meaningful
root and diverges from the other.

\paragraph{Layer two: differential equations.} Once the machine exists its mass
is constant and the flight dynamics begin (\cref{sec:dynamics}). We derive the
Newton-Euler equations on $\SEthree$, the rotor wrench and its allocation map,
first-order rotor dynamics, the battery energy state, and an
orientation-dependent aerodynamic model for the screen. The screen is what
makes this vehicle different from a camera drone: it is a large flat plate
rigidly coupled to the airframe through a gimbal, so attitude modulates
translational drag and drag applied off the centre of mass feeds back as a
pitching moment.

\paragraph{Layer three: control.} \Cref{sec:control} constructs the cascade
from position to velocity to attitude to body rate to rotor speed, states the
exponential stability result of Lee, Leok and McClamroch for the geometric
attitude law on $\SOthree$ together with the ways the implemented law departs
from it, and treats control allocation under saturation. Two items there are
specific to this vehicle. The control authority of each axis is computed as a
linear program over the rotor speed limits, and the attitude gains are bounded
by it, which matters because a rotor optimised for quietness is a rotor with
very little reaction torque and therefore very little yaw authority.
\Cref{prop:governor} states what the reference governor with its moving
keep-out constraint guarantees and what it does not; rate limiting alone,
introduced to stop the display being whipped around the user's head, will on a
curved path steer the corner-cutting reference through the user.

\paragraph{Coupling and computation.} \Cref{sec:codesign} states the full
problem as a differential-algebraic system with a fixed-point constraint and
gives conditions for the inner mass map to be well posed and for the battery
iteration to contract. \Cref{sec:engine,sec:numerics} describe the engine that
solves it and the numerical methods: a bounded root search for the closure
that reports its own failure, secant-accelerated fixed-point iteration for the
battery loop with the simulator inside it and a verification flight at the
end, fixed-step Runge-Kutta for the trajectory, and differential evolution
over the design vector.

\subsection{Contributions}

The mathematical contributions are \cref{thm:trichotomy,thm:bifurcation,%
thm:iteration,thm:sensitivity,thm:pfold,prop:twall,prop:inner}: a complete
description of the sizing equation including its bifurcation structure and the
convergence behaviour of the natural iteration; exact logarithmic sensitivities
of the critical payload; the condition on the disk-area exponent under which
the fold exists at all, and the separate statement that an endurance ceiling
survives the fold's disappearance at the constant-disk-loading boundary; and
the slope condition under which the component mass closure is uniquely
solvable. The control results (\cref{prop:governor}, the authority
computation and the gain heuristic) are stated with their conditions and are
verified by simulation, not proved.

The modelling contributions are the calibration of the acoustic and motor
submodels against measured hardware (\cref{sec:verification}), and the
near-field corrections of \cref{sec:nearfield}, which close a stated gap and,
more usefully, bound its size.

The engineering contribution is the design study of \cref{sec:results}, whose
main conclusion is that for this class of machine the binding constraint is
acoustic rather than energetic, and that stating a noise figure without stating
whether it is a sound power or a pressure, and whether it is free-field or in a
room, is worth more decibels than any design decision available.

\subsection{What this paper does not do}

No hardware was built. Every number reported is the output of a model whose
assumptions are stated where they are introduced, and \cref{sec:limitations}
collects the places where those assumptions are known to be inadequate. In
particular the rotor aerodynamics are momentum theory with a single profile
term and three near-field corrections; there is no blade-element momentum
theory, no wake model, and no treatment of the recirculating flow that a
machine of this disk loading establishes in a closed room within a minute of
taking off. The acoustic model is a broadband scaling law fitted to four
aircraft, coupled to a single-coefficient Sabine room; tonal content is
computed but not summed into the level.
